\scanpage{80}
\mspage{14}
% archive foliation: 82
paire ne change rien; car si on remplace
$D_i'$ par $\lambda D_i'$, donc $D_i''$ par $\frac1\lambda D_i''$, avec $\lambda\in\mathbb R^{*+}$
voisin de 1, ces deux relations vaudront
pour les disques obtenus dans \inserted{le normalisé} $X''$
(c'est probablement tout à fait faux !)
et on obtiendrait la $=$ surface recollée $X$
mais les $\Gamma_i$ remplacés par des cercles
``voisins'' $\Gamma_i'$.

Le cas-clef à comprendre sans
doute est celui où $(g,\nu)=(0,4)$,
i.e. le mécanisme du recollement
de deux ``pantalons'' de Thurston,
donc pour \deleted{\unsure{ces cas}}\inserted{deux} ens. $\varepsilon',\varepsilon''$ à deux
éléments, fixes, les foncteurs des
\[
\left\{\begin{gathered}
\text{pantalons de type combinatoire }\varepsilon'\amalg\{s\},\quad\varepsilon''\amalg\{t\},\\
\text{recollés suivant }\{s,t\},\quad\deleted{\text{donné}}\ \text{suivant}\\
\text{les données }(\ell_i)_{i\in\varepsilon'\amalg\varepsilon''\amalg\unsure{u}}
\in(\mathbb R^{*+})^{\varepsilon'\amalg\varepsilon''\amalg\unsure{u}},\\
\text{et }\theta\in\mathbb U_{\varepsilon'\unsure{\wedge|\times}\varepsilon''}\quad\text{qui fixe le recollement}
\end{gathered}\right.
\]
vers
